\section*{Hoofdstuk \ref{ch:b2:reaction-free-energy} --- Reactie-entropie en reactie-gibbsenergie}

\begin{solution}{exo:b2:reaction-free-energy:1}
(a) Drie mol gas worden vloeibaar: $\Delta_r S^\circ < 0$. (b) Er ontstaat een gas:
$> 0$. (c) Eén gasmolecuul geeft er twee: $> 0$. (d) Ionen in oplossing
vormen een kristal: $< 0$. (e) Smelten: $> 0$. (f) $\Delta\nu_{\text{gas}} = 0$:
klein, en hier positief (\qty{+20}{J/(K.mol)} volgens de tabellen, omdat \ce{HCl}
minder symmetrisch is dan \ce{H2} en \ce{Cl2}).
% ledger: s0:HCl_g, s0:H2_g, s0:Cl2_g
\end{solution}

\begin{solution}{exo:b2:reaction-free-energy:2}
$\Delta_r S^\circ = 2(192.77) - 191.61 - 3(130.68) = \qty{-198.1}{J/(K.mol)}$;
$\Delta_r G^\circ = -91.9 - 298.15 \times (-0.1981) = \qty{-32.8}{kJ/mol}$:
\omterm{def:b2:reaction-free-energy:exergonic}{exergonisch} bij \qty{298}{K}, al werkt de entropieterm (\qty{+59.1}{kJ/mol})
ertegen.
\end{solution}

\begin{solution}{exo:b2:reaction-free-energy:3}
$\Delta_r S^\circ = 69.95 - 130.68 - \tfrac12(205.15) = \qty{-163.3}{J/(K.mol)}$;
$\Delta_r G^\circ = -285.83 + 298.15 \times 0.1633 = \qty{-237.14}{kJ/mol}$.
Het is de vormingsreactie van vloeibaar water (één mol uit de elementen in
hun referentietoestand), dus haar \omterm{def:b2:reaction-free-energy:reaction-gibbs}{standaardreactie-gibbsenergie} is per
definitie $\Delta_f G^\circ(\ce{H2O}, \text{l})$; de tabellen geven dezelfde
waarde.
\end{solution}

\begin{solution}{exo:b2:reaction-free-energy:4}
Ongeveer \qty{42}{J/(K.mol)} bij \qty{298}{K} en \qty{87}{J/(K.mol)} bij
\qty{1000}{K} (vloeistof). Sprongen: \qty{10.6}{J/(K.mol)} bij \qty{692.7}{K} en
\qty{97.7}{J/(K.mol)} bij \qty{1180.2}{K}. $\Delta_{\text{fus}}H = 10.6 \times
692.7 = \qty{7.3}{kJ/mol}$ en $\Delta_{\text{vap}}H = 97.7 \times 1180.2 =
\qty{115}{kJ/mol}$.
\end{solution}

\begin{solution}{exo:b2:reaction-free-energy:5}
$\Delta_r H^\circ = \qty{44.00}{kJ/mol}$, $\Delta_r S^\circ = 188.84 - 69.95 =
\qty{118.9}{J/(K.mol)}$, $\Delta_r G^\circ = 44.00 - 298.15 \times 0.1189 =
\qty{+8.55}{kJ/mol}$. Bij \qty{298}{K} is de verandering niet reversibel (vloeibaar
water en damp bij \qty{1}{bar} zijn niet in evenwicht), dus $\Delta S \ne
Q/T$: $\Delta_r H^\circ/T = \qty{147.6}{J/(K.mol)}$. De positieve
$\Delta_r G^\circ$ betekent dat water bij \qty{298}{K} niet verdampt in
een atmosfeer van zuivere damp bij \qty{1}{bar}: de damp condenseert. Beide
zijn in evenwicht waar $\Delta_r G^\circ = 0$: $T = 44\,004/118.9 =
\qty{370}{K}$, op \qty{3}{K} na het normale kookpunt
(\qty{373}{K}): de \omterm{def:b2:reaction-free-energy:ellingham-approximation}{benadering van Ellingham} is goed over \qty{75}{K}.
\end{solution}

\begin{solution}{exo:b2:reaction-free-energy:6}
$T_i = 57\,100/175.7 = \qty{325}{K}$. Daaronder is $\Delta_r G^\circ > 0$ en
wordt het kleurloze \ce{N2O4} begunstigd; afkoelen van de buis verschuift het gas naar
\ce{N2O4} en de bruine kleur van \ce{NO2} verbleekt (het kwantitatieve verband met
de samenstelling volgt in \cref{ch:b2:equilibrium-shifts}).
\end{solution}

\begin{solution}{exo:b2:reaction-free-energy:7}
$\Delta_r S^\circ = -\dd\Delta_r G^\circ/\dd T = -(-4.0 + 12.0)/100 =
\qty{-0.080}{kJ/(K.mol)} = \qty{-80}{J/(K.mol)}$ en $\Delta_r H^\circ =
\Delta_r G^\circ + T\Delta_r S^\circ = -12.0 + 300 \times (-0.080) =
\qty{-36.0}{kJ/mol}$. Controle: $\Delta_r G^\circ/T$ gaat van $-0.0400$ naar
$\qty{-0.0100}{kJ/(K.mol)}$, helling $3.00 \times 10^{-4}$ per kelvin, gelijk aan
$-\Delta_r H^\circ/T^2 = 36.0/346^2 = 3.00 \times 10^{-4}$ bij de meetkundig
gemiddelde temperatuur $\sqrt{300 \times 400} = \qty{346}{K}$.
\end{solution}

\begin{solution}{exo:b2:reaction-free-energy:8}
$\Delta_r S^\circ = 186.25 - 5.74 - 2(130.68) = \qty{-80.8}{J/(K.mol)}$,
$\Delta_r G^\circ = -74.87 + 298.15 \times 0.0808 = \qty{-50.8}{kJ/mol}$,
de getabelleerde $\Delta_f G^\circ$ van methaan. Omdat beide termen hetzelfde
teken hebben, wordt methaan instabiel boven $T_i = 74\,870/80.8 \approx
\qty{926}{K}$ (\omterm{def:b2:reaction-free-energy:ellingham-approximation}{benadering van Ellingham}): heet methaan kraakt tot koolstof en
waterstof, een route naar beide.
\end{solution}

\begin{solution}{exo:b2:reaction-free-energy:9}
Substitutie: \ce{CH3CHBrCH3 + OH- -> CH3CH(OH)CH3 + Br-}; eliminatie:
\ce{CH3CHBrCH3 + OH- -> CH2=CHCH3 + H2O + Br-}. Het verschil van de twee
standaardreactie-gibbsenergieën is $\Delta(\Delta_r G^\circ) = 40 - T \times
0.090$, negatief boven $40\,000/90 \approx \qty{440}{K}$. Eliminatie maakt
één molecuul meer dan substitutie: haar entropie is groter, en de entropieterm
groeit evenredig met $T$.
\end{solution}

\begin{solution}{exo:b2:reaction-free-energy:10}
Elk van de $N_A$ moleculen heeft 2 oriëntaties: $W = 2^{N_A}$, dus $S =
k_B\ln 2^{N_A} = R\ln 2 = \qty{5.76}{J/(K.mol)}$. Het kristal is niet
volmaakt geordend: bij lage temperatuur zitten de moleculen vastgevroren in hun
willekeurige oriëntaties en kunnen ze de ene geordende rangschikking die de derde
hoofdwet veronderstelt, niet bereiken.
\end{solution}

\begin{solution}{exo:b2:reaction-free-energy:11}
Integreer $\dd\Delta_r S^\circ/\dd T = c/T$: $\Delta_r S^\circ(T) = \Delta_r
S^\circ(T_0) + c\ln(T/T_0)$; Kirchhoff geeft $\Delta_r H^\circ(T) = \Delta_r
H^\circ(T_0) + c(T - T_0)$; trek $T\Delta_r S^\circ(T)$ af. Voor kalksteen bij
\qty{1100}{K}: Ellingham $178.32 - 1100 \times 0.16064 = \qty{1.62}{kJ/mol}$;
correctie $c[T - T_0 - T\ln(T/T_0)] = -0.00195 \times (801.85 - 1436.3) =
\qty{+1.24}{kJ/mol}$: $\Delta_r G^\circ = \qty{2.86}{kJ/mol}$. De fout,
\qty{1.2}{kJ/mol}, is minder dan één procent van $\Delta_r H^\circ$, maar dicht bij
de \omterm{def:b2:reaction-free-energy:inversion-temperature}{inversietemperatuur} beslist ze over het teken.
\end{solution}

\begin{solution}{exo:b2:reaction-free-energy:12}
(a) $\Delta_r H^\circ = -763.2 + 944.0 = \qty{+180.8}{kJ/mol}$, $\Delta_r
S^\circ = 353.2 + 205.15 - 50.62 - 2(223.08) = \qty{+61.6}{J/(K.mol)}$;
$\Delta_r G^\circ(1000) = 180.8 - 61.6 = \qty{+119.2}{kJ/mol}$: \omterm{def:b2:reaction-free-energy:exergonic}{endergonisch},
geen chlorering. (b) Voor \ce{C + O2 -> CO2} is $\Delta_r G^\circ(1000) =
-393.5 - 2.9 = \qty{-396.4}{kJ/mol}$. De som, \ce{TiO2(s) + 2Cl2(g) + C(s)
-> TiCl4(g) + CO2(g)}, heeft $\Delta_r G^\circ = \qty{-277.2}{kJ/mol}$: koolstof
verbruikt de dizuurstof die de eerste reactie maakt, en de gekoppelde reactie is
sterk \omterm{def:b2:reaction-free-energy:exergonic}{exergonisch}.
\end{solution}

\begin{solution}{pb:b2:reaction-free-energy:1}
\textbf{1.} \ce{CaCO3(s) -> CaO(s) + CO2(g)}.
\textbf{2.} $\Delta_r H^\circ = -635.09 - 393.51 + 1206.92 =
\qty{+178.3}{kJ/mol}$: \omterm{def:b2:reaction-enthalpy:exothermic}{endotherm}.
\textbf{3.} $\Delta_r S^\circ = 39.75 + 213.79 - 92.9 = \qty{+160.6}{J/(K.mol)}$,
positief omdat er een gas uit een vaste stof ontstaat.
\textbf{4.} $\Delta_r G^\circ = 178.32 - 298.15 \times 0.16064 =
\qty{+130.4}{kJ/mol}$.
\textbf{5.} Ja: voor ontleding is $\Delta_r G < 0$ nodig, en hier is ze
sterk positief; het koolstofdioxide van de lucht, ver onder \qty{1}{bar},
verlaagt $\Delta_r G$, maar met veel minder dan \qty{130}{kJ/mol} (volgend hoofdstuk).
\textbf{6.} $\Delta_r C^\circ_p = 42.80 + 37.13 - 81.88 =
\qty{-1.95}{J/(K.mol)}$, heel klein: $\Delta_r H^\circ$ en $\Delta_r S^\circ$
veranderen nauwelijks met $T$.
\textbf{7.} $\Delta_r G^\circ(T) = 178.32 - 0.16064\,T$ (\unit{kJ/mol}).
\textbf{8.} \qty{+17.7}{kJ/mol} bij \qty{1000}{K}; \qty{-30.5}{kJ/mol} bij
\qty{1300}{K}.
\textbf{9.} $T_i = 178\,320/160.64 = \qty{1110}{K}$.
\textbf{10.} $\Delta_r G^\circ/T = 178.32/T - 0.16064$, met als afgeleide
$-178.32/T^2 = -\Delta_r H^\circ/T^2$.
\textbf{11.} Bij $T_i$ is $c[T_i - T_0 - T_i\ln(T_i/T_0)] = -0.00195 \times
(811.9 - 1459.4) = \qty{+1.26}{kJ/mol}$; $\Delta_r G^\circ$ bereikt nul
$1260/160.6 \approx \qty{8}{K}$ hoger, bij ongeveer \qty{1118}{K}: minder dan
één procent.
\textbf{12.} Bij \qty{1000}{K} is $\Delta_r G > 0$: de ontleding kan niet
verlopen (ongebluste kalk zou koolstofdioxide opnemen). Bij \qty{1300}{K} is $\Delta_r G <
0$: kalksteen ontleedt.
\textbf{13.} $\Delta_r G$ wordt positief: de ongebluste kalk neemt het koolstofdioxide op
en wordt weer carbonaat.
\textbf{14.} $\delta S_{\text{c}} = -\Delta_r G\,\dd\xi/T = 30\,506/1300 =
\qty{23.5}{J/K}$ per mol.
\textbf{15.} Ja: \cref{ch:b2:chemical-potential} zal laten zien dat
$\Delta_r G = \Delta_r G^\circ + RT\ln(p_{\ce{CO2}}/p^\circ)$, verlaagd
wanneer de druk van koolstofdioxide onder \qty{1}{bar} ligt, zodat de
ontleding al onder $T_i$ begint.
\textbf{16.} $n = 10^6/56.08 = \qty{1.783e4}{mol}$.
\textbf{17.} $1.783 \times 10^4 \times 178.32 = \qty{3.18e6}{kJ}$, dus
\qty{3.18}{GJ} per ton.
\textbf{18.} $1.783 \times 10^4 \times 44.01 = \qty{785}{kg}$ \ce{CO2}.
\textbf{19.} Toe te voeren warmte $3.18 \times 10^6/0.50 = \qty{6.36e6}{kJ}$, dus
$6.36 \times 10^6/802.3 = \qty{7.93e3}{mol}$ methaan, \qty{127}{kg}.
\textbf{20.} $7.93 \times 10^3 \times 44.01 = \qty{349}{kg}$ \ce{CO2}.
\textbf{21.} $785 + 349 = \qty{1134}{kg}$ per ton ongebluste kalk, waarvan 69\,\%
uit de kalksteen.
\textbf{22.} Het koolstofdioxide van de kalksteen ligt vast door de stoichiometrie,
waarmee de oven ook verwarmd wordt; alleen het aandeel van de brandstof kan omlaag.
\textbf{23.} Kalksteen ontleedt onder \qty{1}{bar} koolstofdioxide boven
zijn \omterm{def:b2:reaction-free-energy:inversion-temperature}{inversietemperatuur}, $\boldsymbol{T_i \approx \qty{1.11e3}{K}}$
(ongeveer \qty{840}{\celsius}).
\end{solution}
